\section{Corrective Actions and Verification}

\subsection{Containment and CAPA Plan}
\makeactiontable{Corrective and Preventive Actions}{
    Reproduce the failure with instrumented switching-node capture. & Failure Team & 2026-04-11 & In Progress \\
    Review component derating against measured stress envelope. & Design Lead & 2026-04-13 & Open \\
    Update fault-code mapping and service guide for the observed condition. & Firmware Lead & 2026-04-16 & Open \\
}

\subsection{Verification Strategy}
\begin{table}[H]
    \centering
    \small
    \begin{tabular}{p{3.5cm} p{3.1cm} p{1.8cm} p{3.1cm}}
        \toprule
        \textbf{Verification Step} & \textbf{Success Condition} & \textbf{Owner} & \textbf{Comment} \\
        \midrule
        Reproduce under controlled bench setup & Failure signature appears consistently & Lab Team & Confirms the scenario is real and repeatable. \\
        Apply corrective design or control change & Symptom no longer appears & Design Team & Validate across nominal and corner conditions. \\
        Re-run guardband checks & No new regression introduced & Validation Team & Check thermal, startup, and abnormal-load conditions. \\
        \bottomrule
    \end{tabular}
    \caption{Recommended failure-closure verification matrix}
    \label{tab:failure-verification}
\end{table}

\riskbox{
Do not close the investigation just because the symptom disappeared once. The
final section should state what evidence proves the issue is fixed and what
residual risk remains.
}
